\scanpage{154}\mspage{28}
\uline{\unsure{Dém|Dèm}} Tout a été \unsure{prouvé|prouve}\unsure{\}|;|?} si on se rappelle que les objets de $M^r$ sont $r^{\pi_0}$-accessibles (cf. 3.2).

\deleted{\uline{Corollaire} Les propriétés de stabilité \unsure{restent|sont} valables, lorsque $s$ est}

4. Compléments sur les structures de Quillen.

\uline{Définition 4.1} Soit $M$ une catégorie. On appelle \uline{\unsure{précouple|prêcouple} de Quillen} dans $M$ un couple $(\Phi,\Psi)$ de parties de $\mathrm{Fl}\,M$, satisfaisant les deux conditions suivantes

a) On a $\Phi\longleftrightarrow\Psi$, i.e. $\Psi\subset\Phi_*$, ou ce qui revient au $=$, $\Phi\subset\Psi^*$ \unsure{\&|$\epsilon$|;} -- soit encore : Pour tout carré commutatif dans $M$
\diagram{
\begin{tikzcd}[ampersand replacement=\&]
A\arrow[r,dash,"u","{\unsure{\phantom{\longrightarrow}}}"']\arrow[d,dash,"i"',"{\unsure{\phantom{\downarrow}}}"]\&X\arrow[d,"{\unsure{p|r}}"]\\
\unsure{B|P}\arrow[r,dash,"{\unsure{v|\nu}}"']\arrow[ur,dotted,dash,"{\unsure{w|\omega}}","{\unsure{\phantom{\longrightarrow}}}"']\&Y
\end{tikzcd}
\quad(4.1.1.)
}
avec $i\in\Phi$, $p\in\Psi$, il existe $w$ \unsure{telle|tel} que $wi=u$, $pw=v$.
\margin{\unsure{;2|;z|i2}}

b) Tout $f\in\mathrm{Fl}\,M$ se factorise en
\[
f=pi,\qquad\text{avec }i\in\Phi,\ p\in\Psi.
\]

On écrit
\[
\Phi\mathrel{\begin{tikzpicture}[baseline=-0.45ex]
\draw[<->] (-0.65,0)--(0.65,0);
\draw (-0.04,0.14)--(0.06,-0.18);
\draw (-0.36,0.04)--(-0.27,0.09);
\end{tikzpicture}}\Psi
\tag*{(4.1.2)}
\]
pour indiquer que $(\Phi,\Psi)$ est un \inserted{pré}couple de Quillen.
On \deleted{appelle} dit que le couple de Quillen $(\Phi,\Psi)$ est \uline{clos}, si $\Phi$ est clos à gauche, $\Psi$ clos à droite.
On va voir plus bas \inserted{(voir p.31)} que cela implique (donc \unsure{équivaut|revient} à)
\[
\Phi=\Psi^*,\qquad\Psi=\Phi_{\unsure{\alpha|*}}.
\tag*{(4.1.3.)}
\]
% archive sheet number: 153
